% BHU PYQ -- Professional Exam Template
\documentclass[11pt,a4paper]{article}

\usepackage[a4paper,top=2.5cm,bottom=2.5cm,left=2.8cm,right=2.8cm,headheight=14pt]{geometry}
\usepackage{amsmath,amssymb}
\usepackage[T1]{fontenc}
\usepackage[utf8]{inputenc}
\usepackage{lmodern}
\usepackage{microtype}
\usepackage{fancyhdr}
\usepackage{enumitem}
\usepackage{listings}
\usepackage{hyperref}
\usepackage{lastpage}
\usepackage{parskip}

\hypersetup{colorlinks=true,urlcolor=black,linkcolor=black,
  pdftitle={BVCA-403: Cloud Computing -- BHU PYQ}}

\pagestyle{fancy}
\fancyhf{}
\renewcommand{\headrulewidth}{0.4pt}
\renewcommand{\footrulewidth}{0.4pt}
\fancyhead[L]{\small\textit{Banaras Hindu University}}
\fancyhead[R]{\small\textit{BVCA-403: Cloud Computing}}
\fancyfoot[C]{\small Page \thepage\ of \pageref{LastPage}}

\lstset{
  basicstyle=\small\ttfamily,
  frame=single,
  breaklines=true,
  columns=flexible,
  keepspaces=true
}

\newlist{parts}{enumerate}{2}
\setlist[parts,1]{label=(\alph*),leftmargin=*,itemsep=4pt,topsep=3pt}
\setlist[parts,2]{label=(\roman*),leftmargin=*,itemsep=2pt,topsep=2pt}
\newcommand{\pts}[1]{\hfill\textit{[#1]}}

\begin{document}
\begin{center}
  {\large\bfseries BANARAS HINDU UNIVERSITY}\\[2pt]
  {\normalsize B.Voc. Semester IV Examination, 2022-23}\\[6pt]
  \rule{0.8\linewidth}{0.5pt}\\[6pt]
  {\large\bfseries Paper: BVCA-403 --- Cloud Computing}\\[4pt]
  {\normalsize Time: Three Hours \hfill Maximum Marks: 70}
\end{center}

\rule{\linewidth}{0.4pt}

\smallskip
\noindent\textit{Instructions: This question paper carries three sections :}

\bigskip

. REFNO 065163
| Paper : BVCA-403
| Critical Design Thinking in social Science
| Note: This question paper carries three sections :
| Section-A, B and C. All questions are
: compulsory.
| Section -A’
t
| Note: Answer any two of the following each in
| 500 words approximately. Each question.
: carries 12 marks. 12x2=24
| 1. What do you mean by critical thinking ?
Describe the different barriers to critical
thinking.
P.T.O,
REFNO 065163 |

\medskip\noindent\textbf{Question 2.} Explain the meaning of thinking. Differentiate
between critical thinker and uncritical
thinker.

\medskip\noindent\textbf{Question 3.} Whatis design thinking ? Describe the phases
of design thinking. /

\medskip\noindent\textbf{Question 4.} What are essential design thinking skills'?
Discuss any two design thinking skills. *

\bigskip\noindent{\large\bfseries SECTION --- B}
\rule{\linewidth}{0.4pt}\smallskip

Note: Answer any six of the following questions.
in approximately 200 words. Each question
carries 6 marks : 6x6=36

\medskip\noindent\textbf{Question 5.} (a) Describe the benefits of critical thinking!
\begin{parts}
  \item Highlight the importance of critical \_
  thinking in life. |
  \item How critical thinking helps employees at
  4 workplace ? Discuss.
  * (dj) Discuss how critical thinking makes
  personal life better.
  \item Why design thinking skills matter ?
  Discuss.
  : REFNO 065163
  \item Describe how one can learn design
  thinking.
  \item Differentiate between hard skill and soft
skills.
  \item Describe the functions of communication.
  \item Suggest strategies to improve
  communications skills.
\end{parts}

\bigskip\noindent{\large\bfseries SECTION --- C}
\rule{\linewidth}{0.4pt}\smallskip

Note: Answer each of the following questions ina
word or a sentence. Each question carries
1 mark. 10x1=10

\medskip\noindent\textbf{Question 6.} (a) Who introduced the term critical thinking ?-
, | (b) Explain consistency as critical thinking
standard. |
{c) Differentiate between argument and
| explanation. .
(ad) Why people with critical thinking skills are |
in demand by organizations ?
\begin{parts}
  \item Why should we study logic ?
  REFNO 065163
  \item What is meant by reflective thinking ?
  \item What is not an argument ?: .
  \item Give the meaning of argument.
  \item What is meant by analytical skills ?
  (j} How one can practice active listening ?
  kik
\end{parts}

\vfill
\rule{\linewidth}{0.4pt}
\begin{center}\small
\href{https://bhu-exam.vercel.app/}{Click here to visit our website}\\[4pt]\href{https://me-aryan.vercel.app/}{Click here to visit our contributor}\\[4pt]\href{https://quashan-me.vercel.app/}{Click here to visit our contributor}
\end{center}
\end{document}
